% Part: turing-machines
% Chapter: undecidability
% Section: unsolvability-decision-problem

\documentclass[../../../include/open-logic-section]{subfiles}

\begin{document}

\olfileid{tur}{und}{uns}
\olsection{استحالة حل مسألة القرار}

\begin{thm}
\ollabel{thm:decision-prob}
مسألة القرار غير قابلة للحل: لا توجد آلة تورنغ~$D$ تتوقف في نهاية
المطاف إذا بدأت على شريط يحتوي !!a{sentence}~$!B$ من منطق الرتبة
الأولى دخلًا، وتخرج~$1$ إذا وفقط إذا كانت $!B$ صادقة منطقيًا، وتخرج
$0$ خلاف ذلك.
\end{thm}

\begin{proof}
افترض أن مسألة القرار قابلة للحل، أي افترض وجود آلة تورنغ~$D$ تحلها.
عندئذ نستطيع حل مسألة التوقف كما يأتي. ننشئ آلة تورنغ~$E$ تحسب، إذا
أعطيت دخلًا العدد~$e$ لآلة تورنغ~$M_e$ والدخل~$w$، ال!!{sentence}
المقابلة $!T(M_e,w)\lif !E(M_e,w)$، ثم تتوقف وهي تمسح الخانة الواقعة
في أقصى يسار الشريط. وعندئذ تبدأ الآلة $E\concat D$، إذا أعطيت
الدخلين $e$ و$w$، بحساب~$!T(M_e,w)\lif !E(M_e,w)$، ثم تشغل آلة
مسألة القرار~$D$ على ذلك الدخل. وتتوقف $D$ بخرج~$1$ إذا وفقط إذا
كانت $!T(M_e,w)\lif !E(M_e,w)$ صادقة منطقيًا، وتخرج~$0$ خلاف ذلك.
وتفيد \olref[ver]{lem:halt-if-valid} و
\olref[ver]{lem:valid-if-halt} بأن $!T(M_e,w)\lif !E(M_e,w)$ صادقة
منطقيًا إذا وفقط إذا توقفت $M_e$ عند الدخل~$w$. ومن ثم تتوقف
$E\concat D$، عند الدخلين $e$ و$w$، بخرج~$1$ إذا وفقط إذا توقفت
$M_e$ عند الدخل~$w$، وتتوقف بخرج~$0$ خلاف ذلك. أي إن
$E\concat D$ ستحل مسألة التوقف. لكننا نعلم من
\olref[hal]{thm:halting-problem} استحالة وجود آلة تورنغ كهذه.
\end{proof}

\begin{cor}\ollabel{cor:undecidable-sat}%
ليس قابلًا للقرار ما إذا كانت !!{sentence} اعتباطية من منطق الرتبة
الأولى قابلة للإرضاء.
\end{cor}

\begin{proof}
  افترض أن قابلية الإرضاء تقررها آلة تورنغ~$S$. عندئذ يمكننا حل
  مسألة القرار كما يأتي: إذا أعطيت !!a{sentence}~$B$ دخلًا، فانقل
  $!B$ خانة واحدة نحو اليمين. ثم عد إلى الخانة~$1$ واكتب الرمز
  $\lnot$.

  شغل الآن آلة تورنغ~$S$. وستتوقف في نهاية المطاف وعلى الشريط خرج
  إما $1$ (إذا كانت $\lnot !B$ قابلة للإرضاء) وإما~$0$ (إذا كانت
  $\lnot !B$ غير قابلة للإرضاء). فإذا كانت ثمة~$\TMstroke$ في
  الخانة~$1$ فامحها؛ وإذا كانت الخانة~$1$ خالية فاكتب~$\TMstroke$،
  ثم توقف.

  تتوقف آلة تورنغ هذه دائمًا، ويكون خرجها~$1$ إذا وفقط إذا كانت
  $\lnot !B$ غير قابلة للإرضاء، و$0$ خلاف ذلك. ولما كانت $!B$ صادقة
  منطقيًا إذا وفقط إذا كانت $\lnot !B$ غير قابلة للإرضاء، فالآلة
  تخرج~$1$ إذا وفقط إذا كانت $!B$ صادقة منطقيًا، وتخرج~$0$ خلاف
  ذلك؛ أي إنها ستحل مسألة القرار.
\end{proof}

\begin{explain}
إذن لا توجد آلة تورنغ تعطي دائمًا جوابًا صحيحًا بـ«نعم» أو «لا» عن
السؤال «هل $!B$ !!{sentence} صادقة منطقيًا في منطق الرتبة الأولى؟».
لكن \emph{توجد} آلة تورنغ تعطي دائمًا جواب «نعم» صحيحًا، ولا تتوقف
ببساطة إذا كان الجواب «لا». وينتج ذلك من مبرهنتي السلامة والتمام
لمنطق الرتبة الأولى، ومن إمكان تعداد ال!!{derivation}s حسابيًا.
\end{explain}

\begin{thm}
  \ollabel{thm:valid-ce}%
  الصدق المنطقي ل!!{sentence}s من الرتبة الأولى شبه قابل للقرار: توجد
  آلة تورنغ~$E$ تتوقف في نهاية المطاف وتخرج~$1$ إذا وفقط إذا كانت
  $!B$ صادقة منطقيًا، حين تبدأ على شريط يحتوي !!a{sentence}~$!B$ من
  منطق الرتبة الأولى دخلًا؛ ولا تتوقف خلاف ذلك.
\end{thm}

\begin{proof}
  يمكن توليد جميع ال!!{derivation}s الممكنة في منطق الرتبة الأولى،
  واحدًا بعد آخر، بخوارزمية فعالة. وتفعل الآلة~$E$ ذلك، فإذا وجدت
  !!a{derivation} يبين أن $\Proves !B$ توقفت بخرج~$1$. وبمبرهنة
  السلامة، إذا توقفت $E$ بخرج~$1$ فلأن~$\Entails !B$. وبمبرهنة
  التمام، إذا كان $\Entails !B$ وجد !!a{derivation} يبين أن
  $\Proves !B$. ولما كانت $E$ تولد منهجيًا جميع ال!!{derivation}s
  الممكنة، فستجد في نهاية المطاف اشتقاقًا يبين~$\Proves !B$، ومن ثم
  ستتوقف في النهاية بخرج~$1$.
\end{proof}

\end{document}
